\documentclass{article}
\usepackage{graphicx}
\usepackage{amsmath,amssymb,amsthm}
\usepackage{booktabs}
\usepackage{array}

\newtheorem{proposition}{Proposition}

\begin{document}

% =============================================================================
% Case 3 Detailed Explanation — Unstable ReLU Encoding
% Corresponds to: _hz_apply_relu() in tf_mlp.py, lines 209--314
% =============================================================================

\subsection{Case 3: Unstable Neurons --- Detailed Step-by-Step}
\label{sec:relu_case3_detailed}

\subsubsection{Goal}

For an unstable neuron $h$ with bounds $\alpha = \underline{z}_h < 0 < \overline{z}_h = \beta$, we want to \emph{exactly} encode:
\begin{equation}
y_h = \mathrm{ReLU}(x_h) = \begin{cases} x_h & \text{if } x_h \geq 0, \\ 0 & \text{if } x_h < 0. \end{cases}
\end{equation}

The difficulty: ReLU is piecewise linear with two branches switching at $x = 0$. The hybrid zonotope must use a binary generator to select the branch exactly, without introducing any over-approximation.

% -----------------------------------------------------------------------
\subsubsection{Step 1: Introduce New Variables}
\label{sec:step1_vars}

% Code: lines 239--265, out_Gc, out_Gb, out_c, col_xi1..col_z

We introduce five new generators per unstable neuron:

\begin{center}
\begin{tabular}{@{}llll@{}}
\toprule
Variable & Type & Range & Purpose \\
\midrule
$\xi_1$ & continuous & $[-1, 1]$ & Parameterizes $x_h$ on the inactive branch ($x \leq 0$) \\
$\xi_2$ & continuous & $[-1, 1]$ & Parameterizes $y_h$ on the active branch ($x \geq 0$) \\
$\xi_3$ & continuous & $[-1, 1]$ & Slack variable for $\xi_1$ (absorbs residual in equality) \\
$\xi_4$ & continuous & $[-1, 1]$ & Slack variable for $\xi_2$ (absorbs residual in equality) \\
$z$     & binary     & $\{-1, +1\}$ & Branch selector: $+1 \to$ inactive, $-1 \to$ active \\
\bottomrule
\end{tabular}
\end{center}

\paragraph{Why four continuous generators, not two?}
Variables $\xi_3$ and $\xi_4$ are \emph{slack} variables. They do not appear in the output formula (their rows in $G^c_{\mathrm{out}}$ are zero for neuron $h$). Their sole purpose is to absorb the residual in the equality constraints, enabling us to use \textbf{pure equalities} instead of inequalities. Without slack variables, encoding ``$\xi_2$ is forced to 1 when $z = +1$'' would require the \emph{inequality} $|\xi_2| \leq (1 - z)/2$, violating the paper's Definition~3.1 which mandates $A^c \xi^c + A^b \xi^b = b$.

In the constraint matrices, $\xi_1, \xi_2, \xi_3, \xi_4$ occupy four new columns in $G^c_{\mathrm{out}}$ (columns $n_g, n_g{+}1, n_g{+}2, n_g{+}3$ for the $j$-th unstable neuron), and $z$ occupies one new column in $G^b_{\mathrm{out}}$ (column $n_b{+}j$).

% -----------------------------------------------------------------------
\subsubsection{Step 2: Output Encoding}
\label{sec:step2_output}

% Code: line 268, out_c[unstable] = beta/2
% Code: line 269, out_Gc[unstable, col_xi2] = -beta/2

The output center and generators for neuron $h$ are set as:
\begin{equation}
\label{eq:case3_output}
\boxed{y_h = \frac{\beta}{2} + \left(-\frac{\beta}{2}\right) \cdot \xi_2 = \frac{\beta}{2}(1 - \xi_2).}
\end{equation}

In terms of the HZ data structure:
\begin{itemize}
    \item $c_{\mathrm{out}}[h] = \beta / 2$,
    \item $G^c_{\mathrm{out}}[h, \mathrm{col}_{\xi_2}] = -\beta / 2$,
    \item all other new generator entries for row $h$ are zero.
\end{itemize}

\paragraph{Why this form?}
\begin{itemize}
    \item When $\xi_2 = +1$: $y_h = \tfrac{\beta}{2}(1-1) = 0$ \quad (desired output when inactive).
    \item When $\xi_2 = -1$: $y_h = \tfrac{\beta}{2}(1+1) = \beta$ \quad (maximum output when active).
    \item For $\xi_2 \in [-1,1]$: $y_h \in [0, \beta]$ continuously.
\end{itemize}

The key question is: \emph{how do we force $\xi_2 = 1$ when the neuron is inactive, and let $\xi_2$ be free when active?} The answer lies in the graph equalities (Step~3).

% -----------------------------------------------------------------------
\subsubsection{Step 3: Graph Equalities}
\label{sec:step3_graph}

% Code: lines 283--293 (eq_Ac, eq_Ab, eq_b for eq1, eq2)

Two equality constraints encode the binary switching structure:
\begin{align}
\xi_1 + \xi_3 + z &= 1, \label{eq:case3_graph1} \\
\xi_2 + \xi_4 - z &= 1. \label{eq:case3_graph2}
\end{align}

These are stored as rows in the constraint matrices. Specifically, for the $j$-th unstable neuron (row indices $3j$ and $3j{+}1$ in $A^c_{\mathrm{eq}}$):

\begin{equation}
\label{eq:graph_eq_matrix}
\underbrace{
\begin{pmatrix}
\cdots & 1_{(\xi_1)} & 0 & 1_{(\xi_3)} & 0 & \cdots \\
\cdots & 0 & 1_{(\xi_2)} & 0 & 1_{(\xi_4)} & \cdots
\end{pmatrix}
}_{A^c_{\mathrm{eq}}[\text{rows } 3j,\, 3j{+}1]}
\boldsymbol{\xi}^c
+
\underbrace{
\begin{pmatrix}
\cdots & 1_{(z)} & \cdots \\
\cdots & -1_{(z)} & \cdots
\end{pmatrix}
}_{A^b_{\mathrm{eq}}[\text{rows } 3j,\, 3j{+}1]}
\boldsymbol{\xi}^b
=
\underbrace{
\begin{pmatrix}
1 \\ 1
\end{pmatrix}
}_{b_{\mathrm{eq}}[\text{rows } 3j,\, 3j{+}1]}.
\end{equation}

\paragraph{Case analysis: $z = +1$ (inactive branch, $x_h \leq 0$).}

From \eqref{eq:case3_graph1}:
\[
\xi_1 + \xi_3 + 1 = 1 \implies \xi_1 + \xi_3 = 0 \implies \xi_3 = -\xi_1.
\]
Since $\xi_1 \in [-1,1]$ and $\xi_3 = -\xi_1 \in [-1,1]$, this is always feasible. \textbf{$\xi_1$ is free.}

From \eqref{eq:case3_graph2}:
\[
\xi_2 + \xi_4 - 1 = 1 \implies \xi_2 + \xi_4 = 2.
\]
Since $\xi_2, \xi_4 \in [-1,1]$, the maximum of $\xi_2 + \xi_4$ is $1 + 1 = 2$. This is achievable \emph{only} when $\xi_2 = 1$ and $\xi_4 = 1$. \textbf{$\xi_2 = 1$ is forced.}

Therefore: $y_h = \tfrac{\beta}{2}(1 - 1) = 0$. \checkmark

\paragraph{Case analysis: $z = -1$ (active branch, $x_h \geq 0$).}

From \eqref{eq:case3_graph1}:
\[
\xi_1 + \xi_3 - 1 = 1 \implies \xi_1 + \xi_3 = 2 \implies \xi_1 = \xi_3 = 1.
\]
\textbf{$\xi_1 = 1$ is forced.}

From \eqref{eq:case3_graph2}:
\[
\xi_2 + \xi_4 + 1 = 1 \implies \xi_2 + \xi_4 = 0 \implies \xi_4 = -\xi_2.
\]
\textbf{$\xi_2$ is free} $\in [-1, 1]$, giving $y_h = \tfrac{\beta}{2}(1 - \xi_2) \in [0, \beta]$.

\paragraph{Summary of the switching mechanism.}

\begin{center}
\begin{tabular}{@{}ccccc@{}}
\toprule
$z$ & $\xi_1$ & $\xi_2$ & $y_h$ & Branch \\
\midrule
$+1$ & free $\in [-1,1]$ & \textbf{forced} $= 1$ & $0$ & inactive ($x \leq 0$) \\
$-1$ & \textbf{forced} $= 1$ & free $\in [-1,1]$ & $\frac{\beta}{2}(1-\xi_2) \in [0,\beta]$ & active ($x \geq 0$) \\
\bottomrule
\end{tabular}
\end{center}

The graph equalities implement a ``binary switch'': the value of $z$ determines which variable is locked and which is free. The slack variables $\xi_3, \xi_4$ absorb the residual, enabling this mechanism through pure equalities.

% -----------------------------------------------------------------------
\subsubsection{Step 4: Linking Equality}
\label{sec:step4_linking}

% Code: lines 296--303

At this point, $y_h$ on the active branch takes values in $[0, \beta]$, but we have not yet enforced $y_h = x_h$ (as opposed to an arbitrary positive value). The \textbf{linking equality} provides this crucial connection.

The original input for neuron $h$ is expressed through the input HZ generators:
\begin{equation}
x_h = c_h + G^c[h,:]\,\boldsymbol{\xi}^c_{\mathrm{old}} + G^b[h,:]\,\boldsymbol{\xi}^b_{\mathrm{old}}.
\end{equation}

The linking equality (row $3j{+}2$ in $A^c_{\mathrm{eq}}$) is:
\begin{equation}
\label{eq:case3_linking}
\boxed{
\frac{\alpha}{2}\,\xi_1 - \frac{\beta}{2}\,\xi_2 + \frac{\alpha}{2}\,z - G^c[h,:]\,\boldsymbol{\xi}^c_{\mathrm{old}} - G^b[h,:]\,\boldsymbol{\xi}^b_{\mathrm{old}} = c_h - \frac{\beta}{2}.
}
\end{equation}

In terms of the constraint matrices (for the $j$-th unstable neuron, row $3j{+}2$):
\begin{equation}
\label{eq:linking_matrix}
A^c_{\mathrm{eq}}[3j{+}2, :] = \bigl[\underbrace{-G^c[h,:]}_{\text{old columns}} \;\big|\; \underbrace{\tfrac{\alpha}{2}}_{(\xi_1)} \;\; \underbrace{-\tfrac{\beta}{2}}_{(\xi_2)} \;\; \underbrace{0}_{(\xi_3)} \;\; \underbrace{0}_{(\xi_4)}\bigr],
\end{equation}
\begin{equation}
A^b_{\mathrm{eq}}[3j{+}2, :] = \bigl[\underbrace{-G^b[h,:]}_{\text{old columns}} \;\big|\; \underbrace{\tfrac{\alpha}{2}}_{(z)}\bigr], \qquad
b_{\mathrm{eq}}[3j{+}2] = c_h - \frac{\beta}{2}.
\end{equation}

\paragraph{Rearranging to express $x_h$:}

From \eqref{eq:case3_linking}:
\begin{equation}
x_h = c_h + G^c[h,:]\,\boldsymbol{\xi}^c_{\mathrm{old}} + G^b[h,:]\,\boldsymbol{\xi}^b_{\mathrm{old}} = \frac{\alpha}{2}\,\xi_1 + \frac{\beta}{2}(1 - \xi_2) + \frac{\alpha}{2}\,z.
\end{equation}

\paragraph{Verification on the active branch ($z = -1$, $\xi_1 = 1$):}
\begin{align}
x_h &= \frac{\alpha}{2} \cdot 1 + \frac{\beta}{2}(1 - \xi_2) + \frac{\alpha}{2} \cdot (-1) = \frac{\beta}{2}(1 - \xi_2). \\
y_h &= \frac{\beta}{2}(1 - \xi_2) \quad (\text{from Step 2}).
\end{align}
Therefore $y_h = x_h$. \checkmark \quad \textbf{Exact!}

Range: $\xi_2 \in [-1,1] \implies x_h \in [0, \beta]$. \checkmark

\paragraph{Verification on the inactive branch ($z = +1$, $\xi_2 = 1$):}
\begin{align}
x_h &= \frac{\alpha}{2}\,\xi_1 + \frac{\beta}{2}(1 - 1) + \frac{\alpha}{2} \cdot 1 = \frac{\alpha}{2}(\xi_1 + 1). \\
y_h &= 0 \quad (\text{from Step 2, since } \xi_2 = 1).
\end{align}
Range: $\xi_1 \in [-1,1] \implies x_h \in [0, \alpha]$. Since $\alpha < 0$, this means $x_h \in [\alpha, 0]$. \checkmark

$y_h = 0 = \mathrm{ReLU}(x_h)$ for $x_h \leq 0$. \checkmark

% -----------------------------------------------------------------------
\subsubsection{Step 5: Full Corner Verification}
\label{sec:step5_corners}

We verify all four corner points of the $(\xi_1, \xi_2, z)$ domain:

\begin{center}
\begin{tabular}{@{}cccccccc@{}}
\toprule
$z$ & $\xi_1$ & $\xi_2$ & $x_h$ & $y_h$ & $\mathrm{ReLU}(x_h)$ & Match? \\
\midrule
$+1$ & $-1$ & $1$ (forced) & $\frac{\alpha}{2}(-1+1) = 0$ & $0$ & $0$ & \checkmark \\[2pt]
$+1$ & $+1$ & $1$ (forced) & $\frac{\alpha}{2}(1+1) = \alpha$ & $0$ & $0$ & \checkmark \\[2pt]
$-1$ & $1$ (forced) & $-1$ & $\frac{\beta}{2}(1+1) = \beta$ & $\beta$ & $\beta$ & \checkmark \\[2pt]
$-1$ & $1$ (forced) & $+1$ & $\frac{\beta}{2}(1-1) = 0$ & $0$ & $0$ & \checkmark \\
\bottomrule
\end{tabular}
\end{center}

All corner points match exactly. Intermediate values are covered by the continuous parameterization of $\xi_1$ (on the inactive branch) and $\xi_2$ (on the active branch).

% -----------------------------------------------------------------------
\subsubsection{Step 6: Assembly into the HZ Constraint Matrices}
\label{sec:step6_assembly}

% Code: lines 305--314

The three equalities from Steps 3--4 must be assembled into the output HZ's constraint matrices $A^c_{\mathrm{out}}, A^b_{\mathrm{out}}, b_{\mathrm{out}}$. Let $k$ denote the number of unstable neurons, $n_g$ the number of existing continuous generators, $n_b$ the number of existing binary generators, and $n_c$ the number of existing constraint rows.

\paragraph{Output generator dimensions.}

The output HZ has expanded generator matrices:
\begin{equation}
G^c_{\mathrm{out}} \in \mathbb{R}^{n \times (n_g + 4k)}, \qquad
G^b_{\mathrm{out}} \in \mathbb{R}^{n \times (n_b + k)}.
\end{equation}

The $4k$ new continuous columns are organized as $k$ blocks of $(\xi_1, \xi_2, \xi_3, \xi_4)$, and the $k$ new binary columns are one $z$ per unstable neuron.

\paragraph{Constraint matrix assembly.}

The existing constraints must be extended horizontally (zero-padded for the new generator columns), and the $3k$ new equality rows are appended vertically:

\begin{equation}
\label{eq:case3_Ac_assembly}
A^c_{\mathrm{out}} =
\begin{bmatrix}
A^c_{\mathrm{old}} & \mathbf{0}_{n_c \times 4k} \\[4pt]
\hline
A^c_{\mathrm{eq}}
\end{bmatrix}
\in \mathbb{R}^{(n_c + 3k) \times (n_g + 4k)},
\end{equation}

\begin{equation}
\label{eq:case3_Ab_assembly}
A^b_{\mathrm{out}} =
\begin{bmatrix}
A^b_{\mathrm{old}} & \mathbf{0}_{n_c \times k} \\[4pt]
\hline
A^b_{\mathrm{eq}}
\end{bmatrix}
\in \mathbb{R}^{(n_c + 3k) \times (n_b + k)},
\end{equation}

\begin{equation}
\label{eq:case3_b_assembly}
b_{\mathrm{out}} =
\begin{bmatrix}
b_{\mathrm{old}} \\[4pt]
\hline
b_{\mathrm{eq}}
\end{bmatrix}
\in \mathbb{R}^{n_c + 3k}.
\end{equation}

\paragraph{Why zero-pad the old constraints?}
The existing constraint rows reference only the original generators (columns $0, \ldots, n_g{-}1$ for $A^c$ and $0, \ldots, n_b{-}1$ for $A^b$). The new generators $(\xi_1, \xi_2, \xi_3, \xi_4, z)$ do not participate in any old constraint, so the corresponding entries are zero.

\paragraph{Structure of $A^c_{\mathrm{eq}}$ (new equality rows).}

For each unstable neuron $j \in \{0, \ldots, k{-}1\}$ with neuron index $h_j$, bounds $[\alpha_j, \beta_j]$, and new generator columns $\xi_1^{(j)}, \xi_2^{(j)}, \xi_3^{(j)}, \xi_4^{(j)}$:

\begin{equation}
\label{eq:Ac_eq_detail}
A^c_{\mathrm{eq}} =
\begin{pmatrix}
\text{--- row } 3j \text{ (graph eq 1) ---} \\
\text{--- row } 3j{+}1 \text{ (graph eq 2) ---} \\
\text{--- row } 3j{+}2 \text{ (linking eq) ---}
\end{pmatrix}
\end{equation}

Explicitly, for the $j$-th unstable neuron:

\medskip
\noindent\textbf{Row $3j$ (graph equality 1: $\xi_1 + \xi_3 + z = 1$):}
\[
A^c_{\mathrm{eq}}[3j, :] = [\;\underbrace{0, \ldots, 0}_{n_g} \;|\; \underbrace{0, \ldots, 0, \overbrace{1}^{\xi_1^{(j)}}, 0, \ldots, 0}_{k} \;|\; \underbrace{0, \ldots, 0}_{k} \;|\; \underbrace{0, \ldots, 0, \overbrace{1}^{\xi_3^{(j)}}, 0, \ldots, 0}_{k} \;|\; \underbrace{0, \ldots, 0}_{k}\;]
\]
\[
A^b_{\mathrm{eq}}[3j, :] = [\;\underbrace{0, \ldots, 0}_{n_b} \;|\; \underbrace{0, \ldots, 0, \overbrace{1}^{z^{(j)}}, 0, \ldots, 0}_{k}\;], \qquad
b_{\mathrm{eq}}[3j] = 1.
\]

\medskip
\noindent\textbf{Row $3j{+}1$ (graph equality 2: $\xi_2 + \xi_4 - z = 1$):}
\[
A^c_{\mathrm{eq}}[3j{+}1, :] = [\;\underbrace{0, \ldots, 0}_{n_g} \;|\; \underbrace{0}_{k} \;|\; \underbrace{0, \ldots, 0, \overbrace{1}^{\xi_2^{(j)}}, 0, \ldots, 0}_{k} \;|\; \underbrace{0}_{k} \;|\; \underbrace{0, \ldots, 0, \overbrace{1}^{\xi_4^{(j)}}, 0, \ldots, 0}_{k}\;]
\]
\[
A^b_{\mathrm{eq}}[3j{+}1, :] = [\;\underbrace{0, \ldots, 0}_{n_b} \;|\; \underbrace{0, \ldots, 0, \overbrace{-1}^{z^{(j)}}, 0, \ldots, 0}_{k}\;], \qquad
b_{\mathrm{eq}}[3j{+}1] = 1.
\]

\medskip
\noindent\textbf{Row $3j{+}2$ (linking equality):}
\[
A^c_{\mathrm{eq}}[3j{+}2, :] = [\;\underbrace{-G^c[h_j, :]}_{\text{old columns}} \;|\; \underbrace{\tfrac{\alpha_j}{2}}_{\xi_1^{(j)}} \;\; \underbrace{-\tfrac{\beta_j}{2}}_{\xi_2^{(j)}} \;\; \underbrace{0}_{\xi_3^{(j)}} \;\; \underbrace{0}_{\xi_4^{(j)}} \;|\; \text{other neurons' cols} = 0\;]
\]
\[
A^b_{\mathrm{eq}}[3j{+}2, :] = [\;\underbrace{-G^b[h_j, :]}_{\text{old columns}} \;|\; \underbrace{\tfrac{\alpha_j}{2}}_{z^{(j)}} \;\; \text{other neurons' cols} = 0\;]
\]
\[
b_{\mathrm{eq}}[3j{+}2] = c_{h_j} - \frac{\beta_j}{2}.
\]

\paragraph{Code correspondence.}
\begin{verbatim}
# lines 305--311
old_Ac_ext = cat([hz.Ac, zeros((nc, 4*k))], dim=1)  # (nc, ng+4k)
old_Ab_ext = cat([hz.Ab, zeros((nc,   k))], dim=1)  # (nc, nb+k)

out_Ac = cat([old_Ac_ext, eq_Ac], dim=0)  # (nc+3k, ng+4k)
out_Ab = cat([old_Ab_ext, eq_Ab], dim=0)  # (nc+3k, nb+k)
out_b  = cat([hz.b,       eq_b], dim=0)  # (nc+3k, 1)
\end{verbatim}

\paragraph{Complete output HZ.}

Collecting all components, the output hybrid zonotope is:
\begin{equation}
\label{eq:case3_output_hz}
Z_h^{\mathrm{out}} = \langle\, G^c_{\mathrm{out}},\;\; G^b_{\mathrm{out}},\;\; c_{\mathrm{out}},\;\; A^c_{\mathrm{out}},\;\; A^b_{\mathrm{out}},\;\; b_{\mathrm{out}} \,\rangle
\end{equation}
where:
\begin{align}
G^c_{\mathrm{out}} &\in \mathbb{R}^{n \times (n_g + 4k)}, &
G^b_{\mathrm{out}} &\in \mathbb{R}^{n \times (n_b + k)}, &
c_{\mathrm{out}} &\in \mathbb{R}^{n \times 1}, \\
A^c_{\mathrm{out}} &\in \mathbb{R}^{(n_c + 3k) \times (n_g + 4k)}, &
A^b_{\mathrm{out}} &\in \mathbb{R}^{(n_c + 3k) \times (n_b + k)}, &
b_{\mathrm{out}} &\in \mathbb{R}^{(n_c + 3k) \times 1}.
\end{align}

% -----------------------------------------------------------------------
\subsubsection{Step 7: Why Exact, Not Over-Approximation}
\label{sec:step7_exact}

Classical approaches (DeepPoly, CROWN, triangle relaxation) approximate ReLU by a convex region:
\[
0 \leq y_h \leq \frac{\beta}{\beta - \alpha}(x_h - \alpha),
\]
which allows points where $y_h \neq \mathrm{ReLU}(x_h)$ --- this is an over-approximation.

Our encoding is \textbf{graph-exact}: the linking equality \eqref{eq:case3_linking} \emph{forces} $y_h = x_h$ on the active branch and $y_h = 0$ on the inactive branch. There is no feasible point in the HZ that does not lie on the ReLU graph.

\begin{itemize}
    \item \textbf{For MILP solvers (Gurobi)}: $z \in \{-1, +1\}$ is enforced exactly, so the bounds computed from this HZ are \emph{exact}.
    \item \textbf{For LP solvers (SciPy)}: $z$ is relaxed to $[-1, +1]$, producing an over-approximation of the bounds. However, the HZ \emph{representation itself} remains exact --- only the bounds \emph{computation} introduces relaxation.
\end{itemize}

% -----------------------------------------------------------------------
\subsubsection{Complexity Summary}
\label{sec:complexity}

\begin{center}
\begin{tabular}{@{}lc@{}}
\toprule
Component & Per unstable neuron \\
\midrule
New continuous generators ($\xi_1, \xi_2, \xi_3, \xi_4$) & $+4$ columns in $G^c$ \\
New binary generators ($z$) & $+1$ column in $G^b$ \\
New equality constraints (graph + linking) & $+3$ rows in $(A^c, A^b, b)$ \\
\bottomrule
\end{tabular}
\end{center}

For a layer with $k$ unstable neurons out of $n$ total: the output HZ has $n_g + 4k$ continuous generators, $n_b + k$ binary generators, and $n_c + 3k$ constraint rows.

\end{document}
